% ID: SEL-D0053-OB-E03
% ESTADO: cotejada. Segunda pasada visual y borrador compilado revisados.
% Fuente: Selectividad/sel_2013_matematicas_aplicadas/septiembre/titular/EXAMEN 2-ANDALUCÍA.pdf
% SHA-256: f8633bbd91d33bee02c2cc582c33b7ffeb3643f312b810ad0bb8c5d154d0721d
% Página: 2; opción B; ejercicio original 3.

El 50\% de los préstamos que concede un banco son para vivienda, el 30\% para industria y el 20\% para consumo. No se pagan el 20\% de los préstamos para vivienda, el 15\% de los préstamos para industria y el 70\% de los préstamos para consumo.
\begin{enumerate}[label=\alph*)]
\item \textbf{(1 punto)} Si se elige al azar un préstamo, calcule la probabilidad de que se pague.
\item \textbf{(0.75 puntos)} Se elige un préstamo al azar que resulta impagado, \textquestiondown{}cuál es la probabilidad de que sea un préstamo para consumo?
\item \textbf{(0.75 puntos)} Ante un préstamo impagado el director del banco afirma que es más probable que sea para vivienda que para consumo, \textquestiondown{}lleva razón el director?
\end{enumerate}
